\documentclass[10pt,a4paper]{amsart}
\usepackage{amsmath,amssymb,amsthm,graphicx}
\usepackage{verbatim}
\newenvironment{reference}{\comment}{\endcomment}
%\newenvironment{reference}{}{}
\newenvironment{slogan}{\comment}{\endcomment}
\newenvironment{history}{\comment}{\endcomment}

% For commutative diagrams we use Xy-pic
\usepackage[all]{xy}

% We use 2cell for 2-commutative diagrams.
\xyoption{2cell}
% Theorem environments.
%
\theoremstyle{plain}
\newtheorem{theorem}[subsection]{Theorem}
\newtheorem{proposition}[subsection]{Proposition}
\newtheorem{lemma}[subsection]{Lemma}

\theoremstyle{definition}
\newtheorem{definition}[subsection]{Definition}
\newtheorem{example}[subsection]{Example}
\newtheorem{exercise}[subsection]{Exercise}
\newtheorem{situation}[subsection]{Situation}

\theoremstyle{remark}
\newtheorem{remark}[subsection]{Remark}
\newtheorem{remarks}[subsection]{Remarks}

\numberwithin{equation}{subsection}

% Macros
%
\def\lim{\mathop{\mathrm{lim}}\nolimits}
\def\colim{\mathop{\mathrm{colim}}\nolimits}
\def\Spec{\mathop{\mathrm{Spec}}}
\def\Hom{\mathop{\mathrm{Hom}}\nolimits}
\def\Ext{\mathop{\mathrm{Ext}}\nolimits}
\def\SheafHom{\mathop{\mathcal{H}\!\mathit{om}}\nolimits}
\def\SheafExt{\mathop{\mathcal{E}\!\mathit{xt}}\nolimits}
\def\Sch{\mathit{Sch}}
\def\Mor{\mathop{\mathrm{Mor}}\nolimits}
\def\Ob{\mathop{\mathrm{Ob}}\nolimits}
\def\Sh{\mathop{\mathit{Sh}}\nolimits}
\def\NL{\mathop{N\!L}\nolimits}
\def\CH{\mathop{\mathrm{CH}}\nolimits}
\def\proetale{{pro\text{-}\acute{e}tale}}
\def\etale{{\acute{e}tale}}
\def\QCoh{\mathit{QCoh}}
\def\Ker{\mathop{\mathrm{Ker}}}
\def\Im{\mathop{\mathrm{Im}}}
\def\Coker{\mathop{\mathrm{Coker}}}
\def\Coim{\mathop{\mathrm{Coim}}}

% Boxtimes
%
\DeclareMathSymbol{\boxtimes}{\mathbin}{AMSa}{"02}

%
% Macros for moduli stacks/spaces
%
\def\QCohstack{\mathcal{QC}\!\mathit{oh}}
\def\Cohstack{\mathcal{C}\!\mathit{oh}}
\def\Spacesstack{\mathcal{S}\!\mathit{paces}}
\def\Quotfunctor{\mathrm{Quot}}
\def\Hilbfunctor{\mathrm{Hilb}}
\def\Curvesstack{\mathcal{C}\!\mathit{urves}}
\def\Polarizedstack{\mathcal{P}\!\mathit{olarized}}
\def\Complexesstack{\mathcal{C}\!\mathit{omplexes}}
% \Pic is the operator that assigns to X its picard group, usage \Pic(X)
% \Picardstack_{X/B} denotes the Picard stack of X over B
% \Picardfunctor_{X/B} denotes the Picard functor of X over B
\def\Pic{\mathop{\mathrm{Pic}}\nolimits}
\def\Picardstack{\mathcal{P}\!\mathit{ic}}
\def\Picardfunctor{\mathrm{Pic}}
\def\Deformationcategory{\mathcal{D}\!\mathit{ef}}
\usepackage{cleveref}
\crefname{originalsection}{original section}{original sections}
\Crefname{originalsection}{Original section}{Original sections}
\makeatletter
\newlabel{section-projection}{{0B1N}{source}{}{sourcepointer-1}{}}
\newlabel{section-projection@cref}{{[originalsection][0][]0B1N}{[1][0][]source}{}{}{}}
\makeatother
\crefname{originallemma}{original lemma}{original lemmas}
\Crefname{originallemma}{Original lemma}{Original lemmas}
\makeatletter
\newlabel{lemma-projection-generically-finite}{{0B1P}{source}{}{sourcepointer-2}{}}
\newlabel{lemma-projection-generically-finite@cref}{{[originallemma][0][]0B1P}{[1][0][]source}{}{}{}}
\makeatother
\crefname{originallemma}{original lemma}{original lemmas}
\Crefname{originallemma}{Original lemma}{Original lemmas}
\makeatletter
\newlabel{lemma-projection-generically-immersion}{{0B1Q}{source}{}{sourcepointer-3}{}}
\newlabel{lemma-projection-generically-immersion@cref}{{[originallemma][0][]0B1Q}{[1][0][]source}{}{}{}}
\makeatother
\crefname{originallemma}{original lemma}{original lemmas}
\Crefname{originallemma}{Original lemma}{Original lemmas}
\makeatletter
\newlabel{lemma-projection-generically-immersion-refined}{{0H8A}{source}{}{sourcepointer-4}{}}
\newlabel{lemma-projection-generically-immersion-refined@cref}{{[originallemma][0][]0H8A}{[1][0][]source}{}{}{}}
\makeatother
\crefname{originallemma}{original lemma}{original lemmas}
\Crefname{originallemma}{Original lemma}{Original lemmas}
\makeatletter
\newlabel{lemma-projection-injective}{{0B1R}{source}{}{sourcepointer-5}{}}
\newlabel{lemma-projection-injective@cref}{{[originallemma][0][]0B1R}{[1][0][]source}{}{}{}}
\makeatother
\crefname{originallemma}{original lemma}{original lemmas}
\Crefname{originallemma}{Original lemma}{Original lemmas}
\makeatletter
\newlabel{lemma-transversal-subschemes}{{0B1J}{source}{}{sourcepointer-6}{}}
\newlabel{lemma-transversal-subschemes@cref}{{[originallemma][0][]0B1J}{[1][0][]source}{}{}{}}
\makeatother
\crefname{originalreference}{original reference}{original references}
\Crefname{originalreference}{Original reference}{Original references}
\makeatletter
\newlabel{morphisms-lemma-etale-permanence}{{02GW}{source}{}{sourcepointer-7}{}}
\newlabel{morphisms-lemma-etale-permanence@cref}{{[originalreference][0][]02GW}{[1][0][]source}{}{}{}}
\makeatother
\title{Moving lemma: representatives meeting prescribed subvarieties properly}
\author{The Stacks Project Authors}
\date{stacks-project:intersection.tex:local-snapshot-20261012; GNU FDL 1.2 or later}
\begin{document}
\maketitle
\noindent\textit{Editorial note.} All delivered mathematical text is the author's own, copied verbatim from the pinned Stacks Project chapter file intersection.tex; no statement or proof step was rewritten, reordered or removed.\par\medskip
\noindent\textit{Editorial note.} Span recomputation by label pairing (LANE\_BRIEF 10.2): in this chapter the author's proof environment follows the statement environment as a standalone block; the spec covers the complete environment pair (statement span, proof span) as recorded in the pool candidate.\par\medskip
\noindent\textit{Editorial note.} Occupation scan: combined claim\_key + statement\_anchor + generic-line scan over all registration files (active and withdrawn) found 0 hits for this chapter (is-occupied.json, 2026-10-14); consistent with the duty manager's six-check dispatch.\par\medskip
\ifdefined\hypertarget\hypertarget{sourcepointer-1}{}\fi
\noindent\textit{Source reference.} Original section 0B1N: sources/stacks-chow/tags.txt, source lines 13901--13901. Cross-chapter standard prerequisite cited by the author inside the counted span; it is used as a pointer and its own proof is not part of this item.\par\medskip
\ifdefined\hypertarget\hypertarget{sourcepointer-2}{}\fi
\noindent\textit{Source reference.} Original lemma 0B1P: sources/stacks-chow/tags.txt, source lines 13902--13902. Cross-chapter standard prerequisite cited by the author inside the counted span; it is used as a pointer and its own proof is not part of this item.\par\medskip
\ifdefined\hypertarget\hypertarget{sourcepointer-3}{}\fi
\noindent\textit{Source reference.} Original lemma 0B1Q: sources/stacks-chow/tags.txt, source lines 13903--13903. Cross-chapter standard prerequisite cited by the author inside the counted span; it is used as a pointer and its own proof is not part of this item.\par\medskip
\ifdefined\hypertarget\hypertarget{sourcepointer-4}{}\fi
\noindent\textit{Source reference.} Original lemma 0H8A: sources/stacks-chow/tags.txt, source lines 21485--21485. Cross-chapter standard prerequisite cited by the author inside the counted span; it is used as a pointer and its own proof is not part of this item.\par\medskip
\ifdefined\hypertarget\hypertarget{sourcepointer-5}{}\fi
\noindent\textit{Source reference.} Original lemma 0B1R: sources/stacks-chow/tags.txt, source lines 13904--13904. Cross-chapter standard prerequisite cited by the author inside the counted span; it is used as a pointer and its own proof is not part of this item.\par\medskip
\ifdefined\hypertarget\hypertarget{sourcepointer-6}{}\fi
\noindent\textit{Source reference.} Original lemma 0B1J: sources/stacks-chow/tags.txt, source lines 13897--13897. Cross-chapter standard prerequisite cited by the author inside the counted span; it is used as a pointer and its own proof is not part of this item.\par\medskip
\ifdefined\hypertarget\hypertarget{sourcepointer-7}{}\fi
\noindent\textit{Source reference.} Original reference 02GW: sources/stacks-chow/tags.txt, source lines 3243--3243. Cross-chapter standard prerequisite cited by the author inside the counted span; it is used as a pointer and its own proof is not part of this item.\par\medskip
\begin{lemma}
\label{lemma-moving}
\begin{reference}
See \cite{Roberts}.
\end{reference}
Let $X \subset \mathbf{P}^N$ be a nonsingular closed subvariety.
Let $n = \dim(X)$ and $0 \leq d, d' < n$. Let $Z \subset X$ be a closed
subvariety of dimension $d$ and $T_i \subset X$, $i \in I$ be a finite
collection of closed subvarieties of dimension $d'$. Then there exists
a subvariety $C \subset \mathbf{P}^N$ such that $C$ intersects $X$
properly and such that
$$
C \cdot X = Z + \sum\nolimits_{j \in J} m_j Z_j
$$
where $Z_j \subset X$ are irreducible of dimension $d$, distinct from $Z$, and
$$
\dim(Z_j \cap T_i) \leq \dim(Z \cap T_i)
$$
with strict inequality if $Z$ does not intersect $T_i$ properly in $X$.
\end{lemma}

\begin{proof}
Write $\mathbf{P}^N = \mathbf{P}(V_N)$ so $\dim(V_N) = N + 1$.
We are going to choose a sequence of projections from points
\begin{align*}
& r_N : 
\mathbf{P}(V_N) \setminus \{p_N\} \to \mathbf{P}(V_{N - 1}), \\
& r_{N - 1} :
\mathbf{P}(V_{N - 1}) \setminus \{p_{N - 1}\} \to \mathbf{P}(V_{N - 2}), \\
& \ldots,\\
& r_{n + 1} :
\mathbf{P}(V_{n + 1}) \setminus \{p_{n + 1}\} \to \mathbf{P}(V_n)
\end{align*}
as in Section \ref{section-projection}. At each step we will choose
$p_N, p_{N - 1}, \ldots, p_{n + 1}$ in a suitable Zariski open set.
Pick a closed point $x \in Z \subset X$. For every $i$ pick
closed points $x_{it} \in T_i \cap Z$, at least one in each irreducible
component of $T_i \cap Z$. Taking the composition we obtain
a morphism
$$
\pi = (r_{n + 1} \circ \ldots \circ r_N)|_X :
X \longrightarrow \mathbf{P}(V_n)
$$
which has the following properties
\begin{enumerate}
\item $\pi$ is finite,
\item $\pi$ is \'etale at $x$ and all $x_{it}$,
\item $\pi|_Z : Z \to \pi(Z)$ is an isomorphism
over an open neighbourhood of $\pi(x_{it})$,
\item $T_i \cap \pi^{-1}(\pi(Z)) = (T_i \cap Z) \cup E_i$ with
$E_i \subset T_i$ closed and
$\dim(E_i) \leq d + d' + 1 - (n + 1) = d + d' - n$.
\end{enumerate}
It follows in a straightforward manner from
Lemmas \ref{lemma-projection-generically-finite},
\ref{lemma-projection-generically-immersion},
\ref{lemma-projection-generically-immersion-refined}, and
\ref{lemma-projection-injective} and induction that we can do this;
observe that the last projection is from $\mathbf{P}(V_{n + 1})$ and that
$\dim(V_{n + 1}) = n + 2$ which explains the inequality in (4). In a bit
more detail, let $X_j, Z_j, T_{i, j} \subset \mathbf{P}(V_j)$ be the images
of $X$, $Z$, $T_i$ for $j = N, \ldots, n$. Then
$T_{i, j + 1} \cap r_{j + 1}^{-1}(Z_j) = (T_{i, j + 1} \cap Z_{j + 1})
\cup E_{i, j + 1}$ where we have a dimension bound on $E_{i, j + 1}$
from Lemma \ref{lemma-projection-injective}. Then note that
$T_i \cap \pi^{-1}(\pi(Z))$ is the union of $T_i \cap Z$
and the inverse images of the $E_{i, j}$. Since the transition morphisms
between the $X_j$ are finite we obtain the desired dimension bound.

\medskip\noindent
Let $C \subset \mathbf{P}(V_N)$ be the scheme theoretic closure of
$(r_{n + 1} \circ \ldots \circ r_N)^{-1}(\pi(Z))$. Because $\pi$
is \'etale at the point $x$ of $Z$, we see that the closed subscheme
$C \cap X = \pi^{-1}(\pi(Z))$
contains $Z$ with multiplicity $1$ (local calculation omitted).
Hence by Lemma \ref{lemma-transversal-subschemes} we conclude that
$$
C \cdot X = [Z] + \sum m_j[Z_j]
$$
for some subvarieties $Z_j \subset X$ of dimension $d$. Note that
$$
C \cap X = \pi^{-1}(\pi(Z))
$$
set theoretically. Hence
$T_i \cap Z_j \subset T_i \cap \pi^{-1}(\pi(Z)) \subset (T_i \cap Z) \cup E_i$.
For any irreducible component of $T_i \cap Z_j$ contained in $E_i$ we
have the desired dimension bound. Finally, let $V$ be an irreducible
component of $T_i \cap Z_j$ which is contained in $T_i \cap Z$. To finish
the proof it suffices to show that $V$ does not contain any of the
points $x_{it}$, because then $\dim(V) < \dim(Z \cap T_i)$.
To show this it suffices to show that $x_{it} \not \in Z_j$
for all $i, t, j$.

\medskip\noindent
Set $Z' = \pi(Z)$ and $Z'' = \pi^{-1}(Z')$, scheme theoretically. By
condition (3) we can find an open $U \subset \mathbf{P}(V_n)$ containing
$\pi(x_{it})$ such that $\pi^{-1}(U) \cap Z \to U \cap Z'$ is an isomorphism.
In particular, $Z \to Z'$ is a local isomorphism at $x_{it}$.
On the other hand, $Z'' \to Z'$ is \'etale at $x_{it}$ by condition (2).
Hence the closed immersion $Z \to Z''$ is \'etale at $x_{it}$
(Morphisms, Lemma \ref{morphisms-lemma-etale-permanence}).
Thus $Z = Z''$ in a Zariski neighbourhood of $x_{it}$ which proves
the assertion.
\end{proof}

\end{document}
